\section{Original Cartier orders and residue coefficients}

\begin{theorem}[Actual open restriction preserves the original weighted Cartier order]
\label{thm:intrinsic-numerical-cartier-order-open-pullback}
\uses{thm:fc-curve-cartier-native-local-orders,lem:cartier-principal-equation-regular,thm:fc-stalk-map-original-order-residue,scheme:genuine-ideal-pullback-stalk}
\lean{EnvelopeProperty.NumericalCartierCurveOrders.weightedLocalOrder_comap_open}
\leanok
Let $a:V\to X$ be an actual open immersion of integral locally
Noetherian schemes, let $p:X\to\operatorname{Spec}k$ be original,
and let $I$ be an actual effective Cartier ideal on $X$.
At every original $v\in V$, its weighted native Cartier order at
$a(v)$ equals the weighted native order of the actual comap ideal
$a^{-1}I$ at $v$, over the same structural map $p\circ a$.
No comparison of equations, native orders or residue-field actions is
supplied, and no dimension bound is required.
\end{theorem}
\begin{proof}
Choose the actual Cartier affine equation $c$ of $I$ around $a(v)$,
and choose an actual affine neighbourhood of $v$ inside its inverse
image. The genuine comap ideal on this neighbourhood is generated by
the actual pulled section $c'$. The comap is Cartier since an open
immersion is flat. Thus the original regular-generator theorem makes
$c'$ a nonzerodivisor, in particular nonzero on this nonempty integral
affine neighbourhood. Both chosen equations compute their actual
Cartier orders, independently of the originally chosen Cartier frames.

An open immersion preserves the original point coheight. Outside
coheight one both native orders are zero. At coheight one they are the
lengths of the actual principal stalk quotients. The original
germ-pullback identity identifies the germ of $c'$ with the image of
the germ of $c$ under the same actual stalk map. That map is an
isomorphism, so the two principal orders agree. Its original residue
map is a $k$-algebra isomorphism for the actual structural maps, and
therefore preserves the original residue degree. Multiplying gives
the claimed weighted equality at this same original point.
\end{proof}

\begin{theorem}[The actual affine coefficient algebra has the canonical field degree]
\label{thm:intrinsic-numerical-cartier-affine-coefficient-rank}
\uses{thm:intrinsic-numerical-function-field-degree-affine-rank,scheme:actual-original-flat-global-section-tensor}
\lean{EnvelopeProperty.NumericalCartierCurveOrders.degree_eq_sectionsRank}
\leanok
For an actual finite dominant $f:Y\to\operatorname{Spec}A$ with $A$
an original domain and $Y$ integral, the rank of the whole original
$\Gamma(Y,\mathcal O_Y)$ as an $A$-module, with its actual original
coefficient map, equals the same canonical function-field degree
$N(f)$. Flatness is not required.
\end{theorem}
\begin{proof}
Apply the original affine-rank formula to the whole affine target. The
canonical global-section equivalence
$\Gamma(\operatorname{Spec}A,\mathcal O)\simeq A$ changes only the
coefficient ring. Its inverse followed by the actual $f$-section map is
exactly the original coefficient map to $\Gamma(Y,\mathcal O_Y)$.
The resulting square commutes by the canonical equivalence's two
inverse identities. Rank is preserved by this actual commuting ring
equivalence, giving the formula without a flat base-change theorem.
\end{proof}

\begin{theorem}[The original residue coefficients on Cartier support are finite]
\label{thm:intrinsic-numerical-cartier-support-residue-finite}
\uses{def:fc-affine-scheme-local-coefficients,thm:fc-localized-finite-original-coefficients,thm:fc-curve-cartier-ideal-finite,def:fc-closed-immersion-residue-algebra}
\lean{EnvelopeProperty.NumericalCartierCurveOrders.residueField_finite_of_isFinite,EnvelopeProperty.NumericalCartierCurveOrders.residueField_finite_of_mem_support}
\leanok
Every original point of a finite $k$-scheme has its actual native
residue field finite over $k$, with the original structural action.
Consequently every original point on the support of an effective
Cartier ideal of an integral projective curve has this same property.
No residue extension or coefficient comparison is supplied.
\end{theorem}
\begin{proof}
A finite $k$-scheme is affine, and its actual global algebra is finite
over $k$, hence Artinian. Its localization at the original point's
actual prime is a quotient of that finite algebra, and passage to the
local residue field is another quotient. Both therefore remain finite
over $k$. The canonical original stalk-localization residue algebra
isomorphism identifies this field with the native residue field at
the same original point, preserving the original $k$-action.

The original closed Cartier subscheme of an integral projective curve
is finite over $k$. A point on its actual support is an original point
of that same closed subscheme. Apply the first result there. The
actual closed-immersion residue algebra isomorphism then identifies
its residue field with the original ambient residue field, again over
the same $k$ and the same structural map.
\end{proof}

